%=======================================================================
%  Ejercicio 1 — Un modelo para la tasa de memorización (Tipo 1)
%=======================================================================
\setcounter{ejercicio}{0}   % el enunciado lo numera 1
\begin{ejercicio}[Tipo 1]\label{ej:memorizacion}
\textbf{Un modelo para la tasa de memorización.}
El aprendizaje humano es un proceso extremadamente complejo. Tanto la biología,
como la fisiología, la neurología, la química y la psicología, aun con los
tremendos avances tecnológicos, están lejos de entenderlo. Sin embargo, aunque
esperamos que los modelos más simples de aprendizaje no engloben toda esa
complejidad, nos pueden ayudar a entender aspectos limitados del proceso de
aprendizaje. En este problema vas a estudiar un modelo simple del proceso de
memorización de listas. El modelo está basado en la suposición de que el rango
de aprendizaje es proporcional a la cantidad de lo que queda por aprender. Así
pues, sea $L(t)$ la fracción de la lista ya memorizada al tiempo $t$, de manera
que $L=0$ corresponde a no saber nada de la lista y $L=1$ a haber memorizado ya
la lista entera.

La ecuación diferencial que nos reproduce dicho comportamiento es:
\begin{equation}\label{eq:memorizacion}
  \frac{dL}{dt} = \kappa\,(1-L).
\end{equation}
Todas las personas son diferentes, a distintas personas les tomará distinta
cantidad de tiempo memorizar una lista, es decir, cada quien tiene su valor
personal $\kappa$. La forma de obtener tu valor es la siguiente:
\begin{enumerate}
  \item Usando las listas de números de tres dígitos (\cref{tab:listas})
    deberás recolectar la siguiente información para determinar tu valor de
    $\kappa$ como sigue:
    \begin{enumerate}[label=(\alph*)]
      \item Durante un minuto estudiarás y tratarás de memorizar los elementos
        de la primera lista.
      \item Después pondrás a prueba cuántos números recuerdas escribiendo los
        que recuerdes en el orden correcto respecto a la lista. Si no recuerdas
        un número deja un espacio en blanco.
      \item Durante otro minuto estudiarás la misma lista de números y tratarás
        de memorizarla de nuevo.
      \item Nuevamente, pondrás a prueba tu memoria escribiendo los números que
        recuerdes.
    \end{enumerate}
    Repite este proceso diez veces (o hasta que hayas aprendido la lista por
    completo). Puntúa todas las pruebas que hiciste y realiza una gráfica del
    tiempo $t$, la cantidad de tiempo gastada memorizando la lista, contra $L$,
    la fracción de la lista aprendida.
  \item Usa estos datos para aproximar tu valor $\kappa$ personal y compara los
    datos con las predicciones del modelo. ¿Cuánto tardarás en aprender una
    lista de 20 y 100 números de tres dígitos?
  \item Repite el proceso de la parte 1 con las otras dos listas y calcula tu
    valor de $\kappa$ para esos casos. ¿Es tu valor de $\kappa$ realmente una
    constante o mejora con la práctica? Si suponemos que mejora con la práctica,
    ¿cómo modificarías el modelo para incluir eso?
\end{enumerate}

\begin{table}[H]
  \centering
  \caption{Listas de números de tres dígitos.}
  \label{tab:listas}
  \begin{tabular}{cccc}
    \toprule
    Lista 1 & Lista 2 & Lista 3 & Lista 4 \\
    \midrule
    482 & 145 & 793 & 799 \\
    812 & 909 & 506 & 326 \\
    311 & 969 & 687 & 604 \\
    993 & 206 & 839 & 505 \\
    705 & 532 & 503 & 498 \\
    470 & 507 & 541 & 532 \\
    297 & 232 & 477 & 161 \\
    499 & 536 & 348 & 607 \\
    461 & 281 & 851 & 869 \\
    934 & 534 & 130 & 330 \\
    683 & 391 & 462 & 644 \\
    872 & 332 & 605 & 303 \\
    118 & 270 & 678 & 328 \\
    952 & 442 & 888 & 403 \\
    333 & 683 & 990 & 787 \\
    229 & 812 & 779 & 290 \\
    980 & 718 & 453 & 474 \\
    275 & 960 & 752 & 355 \\
    843 & 145 & 536 & 804 \\
    244 & 721 & 300 & 550 \\
    \bottomrule
  \end{tabular}
\end{table}
\end{ejercicio}

\begin{solucion}
\begin{enumerate}
  \item \pendiente{}
  \item \pendiente{}
  \item \pendiente{}
\end{enumerate}
\end{solucion}
